% =====================================================================
%  STRONA TYTUŁOWA I STRESZCZENIA wg wzoru WEFiM „inz_pl_w12n.docx”
%  (wefim.pwr.edu.pl → Dyplomowanie → Zasady edytowania pracy, p. C)
%  Strona 1: uczelnia, wydział, kierunek, PRACA DYPLOMOWA INŻYNIERSKA,
%            tytuł (identyczny z USOS-APD), autor, opiekun, Wrocław rok.
%  Strona 2: Streszczenie + Słowa kluczowe, Abstract + Keywords
%            (tożsame z tekstem w USOS-APD).
%  Dane wpisuje się w frontmatter/dane_pracy.tex — tu nic nie zmieniać.
% =====================================================================

\newcommand{\stronaTytulowa}{%
  \begin{titlepage}
    \thispagestyle{empty}
    \setstretch{1.0}
    \centering
    {\Large\bfseries Politechnika Wrocławska\par}
    \vspace{2mm}
    {\large Wydział Elektroniki, Fotoniki i~Mikrosystemów\par}
    \vspace{12mm}
    {\large Kierunek studiów: \pracaKierunek\par}
    \ifdefempty{\pracaSpecjalnosc}{}{%
      {\large Specjalność: \pracaSpecjalnosc\par}}
    \vfill
    {\huge\bfseries PRACA DYPLOMOWA\\[2mm] INŻYNIERSKA\par}
    \vspace{14mm}
    {\setstretch{1.15}\LARGE\bfseries \pracaTytul\par}
    \vspace{8mm}
    {\setstretch{1.15}\large\itshape \pracaTytulEN\par}
    \vfill
    \begin{flushright}
      \large
      \begin{tabular}{@{}l@{}}
        Autor:\\[1mm]
        \textbf{\pracaAutor}\\[8mm]
        Opiekun:\\[1mm]
        \textbf{\pracaOpiekun}
      \end{tabular}
    \end{flushright}
    \vfill
    {\large Wrocław \pracaRok\par}
  \end{titlepage}
  % Strona 2 (odwrocie): streszczenia
  \thispagestyle{plain}
  \noindent{\large\bfseries Streszczenie\par}
  \vspace{2mm}
  \pracaStreszczenie\par
  \vspace{3mm}
  \noindent\textbf{Słowa kluczowe:} \pracaSlowaKluczowe\par
  \vspace{10mm}
  \begin{otherlanguage}{english}
  \noindent{\large\bfseries Abstract\par}
  \vspace{2mm}
  \pracaAbstract\par
  \vspace{3mm}
  \noindent\textbf{Keywords:} \pracaKeywords\par
  \end{otherlanguage}
  \cleardoublepage
}
